% GENERATED FILE. Edit canonical lessons, then run npm run generate:books.
% canonical-source-hash: fnv1a64:015bfa225b5d5384
% canonical-lessons: ZH-C152-de-complement, ZH-W152-de-complement, ZH-C152-de-grammar, ZH-C152-liuli, ZH-C152-bucuo, ZH-C152-qingchu, ZH-C152-practice

\chapter{Well Said}
\label{ch:zh-well-said}
\hlchaptermodality{152}

\begin{chapteropening}
\textbf{By the end of this chapter:} \emph{I can say how well someone does something.}

Saying how well something is done with a verb, \zh{得} and an adjective, fluent, not bad and clear, built on one new written character: \zh{得}.
\end{chapteropening}

\section[\texorpdfstring{\zh{得} (de)}{得 (de)}]{\zh{得} (de) — how well it is done}

\label{lesson:ZH-C152-de-complement}



\begin{quote}
A new chapter, and it opens on a small word that says how well.

\begin{itemize}
\raggedright
  \item \emph{Recall:} say \zh{我看到他了} — \emph{wǒ kàndào tā le} — \textbf{R1}, the chapter you have finished
  \item \emph{Recall:} say \emph{zuòwán}, to finish — \textbf{R1}
  \item \emph{Recall:} say \emph{juéde}, to feel; to think — \textbf{R2}, a few chapters back
  \item \emph{Recall:} say \zh{我的书} — \emph{wǒ de shū}, my book — \textbf{R4}, recognition at distance
\end{itemize}
\end{quote}

\subsection*{You'll want to know: \zh{得}}
\textbf{\zh{得}} — \emph{de} — the particle between a verb and how well it is done.

\begin{quote}
\textbf{\zh{你说得很好}} \emph{nǐ shuō de hěn hǎo} — You speak well.
\end{quote}

The next lesson puts the character in your hand. It is the de you already say at the end of \emph{juéde}.

\begin{cousinweb}
\textbf{\zh{得}} (\emph{de}) is written with the character of dé 'to get', inherited from Middle Chinese \emph{tok} (Guangyun, in Baxter's transcription), glossed 'gain and loss'. Its use as a light particle after a verb is a later development, assumed here, not quoted.
\end{cousinweb}

\begin{grammarlens}[title={What you've built}]
\zh{得} and \zh{的} are two different words, and both are said \emph{de}, with no tone.

\zh{的} joins an owner or a description to a noun: \zh{我的书}. \zh{得} joins a verb to how well it went: \zh{说得很好}.
\end{grammarlens}

\subsection*{Your turn}
\begin{itemize}
\raggedright
  \item \emph{Say it:} \emph{de}, light, with no tone
  \item \emph{Say it:} \zh{你说得很好} — \emph{nǐ shuō de hěn hǎo} — You speak well.
  \item \emph{Say it:} \zh{我的书} — \emph{wǒ de shū} — my book: the other de
  \item \emph{Recall:} say \emph{juéde}, then \zh{说得很好}, and hear the same light de
\end{itemize}

\begin{tcolorbox}[breakable,colback=teal!4,colframe=teal!35!black,title={Before you move on}]
What does \zh{得} join? (\textbf{A verb to how well it went}.) And \zh{的}? (\textbf{An owner or a description to a noun}.)
\end{tcolorbox}
\section[\texorpdfstring{\zh{得} (de)}{得 (de)}]{\zh{得} — the left of \zh{很}, then \zh{日}, a bar and \zh{寸}}

\label{lesson:ZH-W152-de-complement}



\begin{quote}
The word first, then the character inside it.

\begin{itemize}
\raggedright
  \item \emph{Recall:} say \emph{de} — \textbf{R1}, one lesson back, before it decays
\end{itemize}

\begin{quote}
\zh{很} \zh{时}
\end{quote}

\emph{Read it:} each one aloud, with nothing said first — \emph{hěn}, \emph{shí}
\end{quote}

\subsection*{Script — the left of \zh{很}, then \zh{日}, a bar and \zh{寸}}
\hlblockfigure{\detokenize{figures/ZH-W152-de-complement-filmstrip.pdf}}{How \zh{得} is written, stroke by stroke}

\begin{quote}
\zh{得}
\end{quote}

\textbf{Four pieces, and you write every one of them already.}

\textbf{\zh{彳}, three strokes}, the left of \zh{很}: a short left-falling stroke, a longer one under it, then the vertical down to the foot.

\textbf{\zh{日}, four strokes}, small, at the upper right: the left side, the top that turns down the right side, the middle bar, the bottom bar.

\textbf{A long bar} under \zh{日}, across the whole right side.

\textbf{\zh{寸}, three strokes}, under the bar, as on the right of \zh{对} and \zh{时}: the bar, the vertical that hooks to the left at its foot, then the dot.

3 + 4 + 1 + 3 = eleven.

\subsection*{Your turn}
\begin{itemize}
\raggedright
  \item \emph{Write it:} \zh{得}, all eleven, with the model in front of you
  \item \emph{Cover:} write it again from memory, and count the pieces, not the strokes
\end{itemize}

\begin{tcolorbox}[breakable,colback=teal!4,colframe=teal!35!black,title={Before you move on}]
How many strokes in \zh{得}? (\textbf{Eleven}.) Which character you write has the same left side? (\textbf{\zh{很}}.) What sits at the foot on the right? (\textbf{\zh{寸}, as in \zh{对} and \zh{时}}.)

Source: \href{https://raw.githubusercontent.com/chanind/hanzi-writer-data/68d10a4b21150cae5e1ebbd223eed289cf32d90c/data/\%E5\%BE\%97.json}{Hanzi Writer Data \zh{得}, pinned stroke paths and medians}.
\end{tcolorbox}
\section[\texorpdfstring{\zh{你说得很好} (nǐ shuō de hěn hǎo)}{你说得很好 (nǐ shuō de hěn hǎo)}]{\zh{你说得很好} — how well, with \zh{得}}

\label{lesson:ZH-C152-de-grammar}



\begin{quote}
Three words in a fixed order, and the middle one is the character you have written.

\begin{itemize}
\raggedright
  \item \emph{Read it:} \zh{说得很好}, aloud — \emph{shuō de hěn hǎo}
  \item \emph{Recall:} say \zh{快} — \emph{kuài}, fast
  \item \emph{Recall:} say \zh{汉语} — \emph{Hànyǔ}, Chinese
\end{itemize}
\end{quote}

\begin{grammarlens}[title={Grammar Lens: the verb, \zh{得}, then how well}]
\begin{quote}
\textbf{After a verb, \zh{得} brings in how well or how fast it is done. The adjective after \zh{得} usually takes \zh{很}.}
\end{quote}

\begin{quote}
\textbf{\zh{你说得很好}} \emph{nǐ shuō de hěn hǎo} — You speak well. \textbf{\zh{他说得很快}} \emph{tā shuō de hěn kuài} — He talks fast.
\end{quote}

With an object, say the verb twice: once with the object, once with \zh{得}:

\begin{quote}
\textbf{\zh{他说汉语说得很快}} \emph{tā shuō Hànyǔ shuō de hěn kuài} — He speaks Chinese fast.
\end{quote}

\zh{很} here is light, as it is in \zh{你很好}: it carries the sentence more than it means \emph{very}.
\end{grammarlens}

\subsection*{Your turn}
Say these aloud:

\begin{itemize}
\raggedright
  \item \zh{你说得很好} — \emph{nǐ shuō de hěn hǎo} — You speak well.
  \item \zh{他说得很快} — \emph{tā shuō de hěn kuài} — He talks fast.
  \item \zh{他说汉语说得很快} — \emph{tā shuō Hànyǔ shuō de hěn kuài} — He speaks Chinese fast.
  \item \zh{你做饭做得很好} — \emph{nǐ zuò fàn zuò de hěn hǎo} — You cook well.
\end{itemize}

\begin{tcolorbox}[breakable,colback=teal!4,colframe=teal!35!black,title={Before you move on}]
What comes after \zh{得}? (\textbf{How well: an adjective, usually with \zh{很}}.) With an object, what do you say twice? (\textbf{The verb: \zh{说汉语说得很快}}.)
\end{tcolorbox}
\section[liúlì]{liúlì — fluent}

\label{lesson:ZH-C152-liuli}



\begin{quote}
Before the new one: read the last lesson's sentence, then a word for slow.

\begin{itemize}
\raggedright
  \item \emph{Read it:} \zh{你说得很好}, aloud — \emph{nǐ shuō de hěn hǎo}
  \item \emph{Recall:} say \emph{màn}, slow — \textbf{R3}, durable
\end{itemize}
\end{quote}

\subsection*{You'll want to know: liúlì}
\emph{liúlì} — \textquotedblleft{}fluent\textquotedblright{}.

\begin{quote}
\emph{tā Hànyǔ shuō de hěn liúlì} — He speaks Chinese fluently.
\end{quote}

Its characters come later; for now it is a word for the ear.

\begin{cousinweb}
\emph{liúlì} is built from liú 'to flow' and lì. Both are inherited from Middle Chinese, \emph{ljuw} and \emph{lijH} (Guangyun, in Baxter's transcription); the Guangyun explains \emph{ljuw}, after an older dictionary, as water moving, and glosses \emph{lijH} 'favourable' and, after the same dictionary, 'sharp'. 'Fluent' is the pair's own sense.
\end{cousinweb}

\begin{grammarlens}[title={What you've built}]
\emph{liúlì} sits after \zh{得} with \zh{很}, as the other adjectives do: \emph{shuō de hěn liúlì}.

Here the topic, Hànyǔ, goes first, so the verb is said only once.
\end{grammarlens}

\subsection*{Your turn}
\begin{itemize}
\raggedright
  \item \emph{Say it:} \emph{liúlì}
  \item \emph{Say it:} \emph{liúlì}, once more
  \item \emph{Say it:} \emph{tā Hànyǔ shuō de hěn liúlì} — He speaks Chinese fluently.
  \item \emph{Recall:} say \emph{màn}, then \emph{liúlì}: slow, then smooth
\end{itemize}

\begin{tcolorbox}[breakable,colback=teal!4,colframe=teal!35!black,title={Before you move on}]
What does \emph{liúlì} mean? (\textquotedblleft{}Fluent\textquotedblright{}.) Where does it go? (\textbf{After \zh{得}, with \zh{很}}.)
\end{tcolorbox}
\section[búcuò]{búcuò — not bad; pretty good}

\label{lesson:ZH-C152-bucuo}



\begin{quote}
Before the new one: last lesson's word, then the word this one says no to.

\begin{itemize}
\raggedright
  \item \emph{Recall:} say \emph{liúlì}, fluent — \textbf{R1}, one lesson back
  \item \emph{Recall:} say \emph{cuò}, wrong — \textbf{R2}, a few chapters back
  \item \emph{Recall:} say \emph{róngyì}, easy — \textbf{R2}
\end{itemize}
\end{quote}

\subsection*{You'll want to know: búcuò}
\emph{búcuò} — \textquotedblleft{}not bad; pretty good\textquotedblright{}.

\begin{quote}
\emph{nǐ zuò de búcuò} — You did pretty well.
\end{quote}

\zh{不} rises before the fall of cuò: \emph{búcuò}, as in \emph{bú qù}.

Its second character comes later; for now it is a word for the ear.

\begin{cousinweb}
\emph{búcuò} is built from \zh{不} 'not' and \emph{cuò} 'wrong', the word you know: not wrong, so pretty good. Both are inherited from Middle Chinese, \emph{pjut} and \emph{tshak} (Guangyun, in Baxter's transcription); the Guangyun glosses \emph{tshak} 'mixed' and 'crossing slantwise', and how it came to mean wrong is not settled.
\end{cousinweb}

\begin{grammarlens}[title={What you've built}]
\emph{búcuò} is praise that sounds modest. After \zh{得} it needs no \zh{很}: \emph{zuò de búcuò}.

On its own it answers how something was: \emph{hěn búcuò}, really quite good.
\end{grammarlens}

\subsection*{Your turn}
\begin{itemize}
\raggedright
  \item \emph{Say it:} \emph{búcuò}
  \item \emph{Say it:} \emph{búcuò}, once more
  \item \emph{Say it:} \emph{nǐ zuò de búcuò} — You did pretty well.
  \item \emph{Recall:} say \emph{cuò}, then \emph{búcuò}
\end{itemize}

\begin{tcolorbox}[breakable,colback=teal!4,colframe=teal!35!black,title={Before you move on}]
What does \emph{búcuò} mean? (\textquotedblleft{}Not bad; pretty good\textquotedblright{}.) What does \zh{不} do before cuò? (\textbf{It rises: \emph{bú}}.)
\end{tcolorbox}
\section[qīngchu]{qīngchu — clear; clearly}

\label{lesson:ZH-C152-qingchu}



\begin{quote}
Before the new one: last lesson's word, then two words for how a thing is.

\begin{itemize}
\raggedright
  \item \emph{Recall:} say \emph{búcuò}, pretty good — \textbf{R1}, one lesson back
  \item \emph{Recall:} say \zh{重要} — \emph{zhòngyào}, important — \textbf{R2}, a few chapters back
  \item \emph{Recall:} say \emph{jiǎndān}, simple — \textbf{R2}
\end{itemize}
\end{quote}

\subsection*{You'll want to know: qīngchu}
\emph{qīngchu} — \textquotedblleft{}clear; clearly\textquotedblright{}.

\begin{quote}
\emph{tā shuō de hěn qīngchu} — He speaks clearly.
\end{quote}

Its characters come later; for now it is a word for the ear.

\begin{cousinweb}
\emph{qīngchu} is built from qīng 'clear' and chu. Both are inherited from Middle Chinese, \emph{tshjeng} and \emph{tsrhjoX} (Guangyun, in Baxter's transcription); the first is glossed 'clear, far from mud'. The Guangyun gives \emph{tsrhjoX} as a plant and a place name, so how it came into 'clear' is not settled.
\end{cousinweb}

\begin{grammarlens}[title={What you've built}]
\emph{qīngchu} answers how well something came across: \emph{shuō de hěn qīngchu}, said clearly.

It is light at the end, \emph{chu} with no tone.
\end{grammarlens}

\subsection*{Your turn}
\begin{itemize}
\raggedright
  \item \emph{Say it:} \emph{qīngchu}
  \item \emph{Say it:} \emph{qīngchu}, once more
  \item \emph{Say it:} \emph{tā shuō de hěn qīngchu} — He speaks clearly.
  \item \emph{Recall:} say \emph{jiǎndān}, then \emph{qīngchu}
\end{itemize}

\begin{tcolorbox}[breakable,colback=teal!4,colframe=teal!35!black,title={Before you move on}]
What does \emph{qīngchu} mean? (\textquotedblleft{}Clear; clearly\textquotedblright{}.) Say he speaks clearly. (\emph{tā shuō de hěn qīngchu}.)
\end{tcolorbox}
\section[Practice]{\zh{他说汉语说得很快}}

\label{lesson:ZH-C152-practice}



\begin{quote}
The chapter closes here, with two people trading compliments.

\begin{itemize}
\raggedright
  \item \emph{Recall:} say \emph{qīngchu}, clear — \textbf{R1}, one lesson back
  \item \emph{Recall:} say \emph{búcuò}, pretty good — \textbf{R1}
  \item \emph{Recall:} say \emph{liúlì}, fluent — \textbf{R2}
\end{itemize}
\end{quote}

\subsection*{The exchange}
\begin{quote}
\textbf{A:} \zh{他说汉语说得很快}
\end{quote}

\begin{quote}
\emph{tā shuō Hànyǔ shuō de hěn kuài} — He speaks Chinese fast.
\end{quote}

\begin{quote}
\textbf{B:} \zh{对} \zh{他说得很好}
\end{quote}

\begin{quote}
\emph{duì, tā shuō de hěn hǎo} — Yes, and he speaks well.
\end{quote}

\begin{quote}
\textbf{A:} \emph{tā Hànyǔ shuō de hěn liúlì} — His Chinese is fluent.
\end{quote}

\begin{quote}
\textbf{B:} \zh{你呢}
\end{quote}

\begin{quote}
\emph{nǐ ne} — And you?
\end{quote}

\begin{quote}
\textbf{A:} \emph{wǒ shuō de hěn màn} — I speak slowly.
\end{quote}

\begin{quote}
\textbf{B:} \emph{nǐ shuō de hěn qīngchu} — You speak clearly.
\end{quote}

\begin{quote}
\textbf{A:} \zh{你做饭做得很好}
\end{quote}

\begin{quote}
\emph{nǐ zuò fàn zuò de hěn hǎo} — You cook well.
\end{quote}

\begin{quote}
\textbf{B:} \emph{búcuò, wǒ juéde hěn jiǎndān} — Not bad. I find it simple.
\end{quote}

Every compliment has the same order: the verb, \zh{得}, then how well. When the object comes along, the verb is said twice; when the language goes first as the topic, once.

\begin{tcolorbox}[breakable,skin=enhanced,colback=orange!6,colframe=orange!45!black,boxrule=0.5pt,arc=1mm,left=6pt,right=6pt,top=4pt,bottom=4pt,fonttitle=\bfseries,title={Read this}]
\emph{Read it:} the exchange through once without stopping, then answer the question below
\end{tcolorbox}

\begin{grammarlens}[title={What you've built — how well}]
The verb, \zh{得}, then how well; say the verb twice when it has an object:

\begin{itemize}
\raggedright
  \item \textbf{\zh{你说得很好}} \emph{nǐ shuō de hěn hǎo} — you speak well.
  \item \textbf{\zh{他说汉语说得很快}} \emph{tā shuō Hànyǔ shuō de hěn kuài} — he speaks Chinese fast.
  \item \emph{tā shuō de hěn qīngchu} — he speaks clearly.
  \item \emph{nǐ zuò de búcuò} — you did pretty well.
\end{itemize}

\textbf{Listen.} \emph{Tā Hànyǔ shuō de zěnmeyàng?}

\textbf{Speak.} Say how well somebody you know cooks, then how fast you speak.

\emph{Write it:} \zh{得} — the left of \zh{很}, then \zh{日}, a long bar, and the right of \zh{对}
\end{grammarlens}

\subsection*{Your turn}
\begin{itemize}
\raggedright
  \item \emph{Say it:} \zh{他说汉语说得很快} — \emph{tā shuō Hànyǔ shuō de hěn kuài} — He speaks Chinese fast.
  \item \emph{Say it:} \emph{nǐ shuō de hěn qīngchu} — You speak clearly.
  \item \emph{Say it:} the whole exchange, both parts
  \item \emph{Write it:} \zh{你说得很好}
\end{itemize}

\begin{tcolorbox}[breakable,colback=teal!4,colframe=teal!35!black,title={Before you move on}]
Say he speaks Chinese fast. (\textbf{\zh{他说汉语说得很快}}.) Say someone speaks fluently. (\emph{tā shuō de hěn liúlì}.) Say not bad. (\emph{búcuò}.)
\end{tcolorbox}
